%% ma: L11-B14-0024
%% lop: 11
%% chuong: 4
%% bai: 14
%% dang: 11.14.1
%% loai: DS
%% muc: [NB, TH, TH, TH]
%% dap_an: "ĐSĐS"
%% nguon: "tai-lieu-hoc-tap-mon-toan-lop-11-theo-chuong-trinh-giao-duc-pho-thong-2018.pdf (Lớp Toán Thầy Duy, Phạm Lê Duy, THPT Chu Văn An), Chương 4 Bài 5, Phần 2, Câu 4 (tr.532)"
%% kien_thuc: "Lớp 11 Bài 14 (phép chiếu song song)"
%% trang_thai: cho-duyet
%% ly_do: "Dạng mới đề xuất (11.14.x), chờ giáo viên duyệt dạng; tài liệu gốc không có lời giải mẫu, đáp án do kho tự giải và kiểm bằng sympy."
\begin{ex}
Cho hình lăng trụ $ABC.A^{\prime}B^{\prime}C^{\prime}$ có $I$ và $I^{\prime}$ lần lượt là trung điểm của đoạn $AB$ và $A^{\prime}B^{\prime}$. Xét tính đúng, sai của các mệnh đề sau
\choiceTF
{\True $AI^{\prime}\parallel IB^{\prime}$.}
{Hình chiếu song song của $I$ trên mặt phẳng $(A^{\prime}B^{\prime}C^{\prime})$ phương $A^{\prime}I$ là điểm $C^{\prime}$.}
{\True Trong $(A^{\prime}B^{\prime}C^{\prime})$, vẽ hình bình hành $A^{\prime}C^{\prime}MI^{\prime}$. Suy ra $ACMI^{\prime}$ là hình bình hành.}
{$\triangle MAA^{\prime}$ là hình chiếu song song của $\triangle CAA^{\prime}$ theo phương $AI^{\prime}$ trên $(A^{\prime}B^{\prime}C^{\prime})$.}
\loigiai{\sloppy a) $AI\parallel I^{\prime}B^{\prime}$ và $AI=I^{\prime}B^{\prime}$ nên $AIB^{\prime}I^{\prime}$ là hình bình hành, suy ra $AI^{\prime}\parallel IB^{\prime}$. Đúng.
b) Đường thẳng qua $I$ song song $A^{\prime}I$ chính là $A^{\prime}I$, cắt $(A^{\prime}B^{\prime}C^{\prime})$ tại $A^{\prime}$. Hình chiếu của $I$ là $A^{\prime}$, không phải $C^{\prime}$. Sai.
c) $\overrightarrow{I^{\prime}M}=\overrightarrow{A^{\prime}C^{\prime}}=\overrightarrow{AC}$ nên $ACMI^{\prime}$ là hình bình hành. Đúng.
d) Theo phương $AI^{\prime}$, hình chiếu của $A$ là $I^{\prime}$ (vì $I^{\prime}\in(A^{\prime}B^{\prime}C^{\prime})$ và nằm trên $AI^{\prime}$), không phải $A^{\prime}$; hình chiếu của $\triangle CAA^{\prime}$ là tam giác $MI^{\prime}A^{\prime}$ chứ không phải $\triangle MAA^{\prime}$. Sai.}
\end{ex}
